\documentclass[
    language=english,
    doctype=manual,
    institution=none,
    titlestyle=book,
    oneside,
]{../../omnilatex}

\RequirePackage{config/document-settings}

\addbibresource{../../bib/bibliography.bib}

\begin{document}

\maketitle

\tableofcontents

%=============================================================================
\chapter{Prerequisites}
\label{sec:prerequisites}
%=============================================================================

Before deploying DataPipe, ensure all infrastructure and software
dependencies are in place. This chapter covers the required hardware,
software, and access permissions.

\section{Infrastructure Requirements}

\begin{table}[htbp]
    \centering
    \caption{Infrastructure requirements for production deployment.}
    \label{tab:infra}
    \begin{tabular}{@{}lll@{}}
        \toprule
        Component & Minimum & Recommended \\
        \midrule
        Application nodes & 3 $\times$ 8 vCPU, 32\,GB RAM &
            3 $\times$ 16 vCPU, 64\,GB RAM \\
        Kafka cluster & 3 $\times$ 4 vCPU, 16\,GB RAM, 500\,GB SSD &
            3 $\times$ 8 vCPU, 32\,GB, 1\,TB NVMe \\
        PostgreSQL & 2 $\times$ 4 vCPU, 16\,GB RAM &
            2 $\times$ 8 vCPU, 32\,GB RAM \\
        Load balancer & 2 vCPU, 4\,GB RAM &
            Managed LB (AWS ALB / GCP CLB) \\
        \bottomrule
    \end{tabular}
\end{table}

\section{Software Dependencies}

\begin{itemize}
    \item Docker 24+ and Docker Compose v2.
    \item Kubernetes 1.28+ (for K8s deployment path).
    \item Helm 3.14+ (for chart-based deployment).
    \item Terraform 1.7+ (for infrastructure provisioning).
    \item AWS CLI or gcloud CLI (depending on cloud provider).
\end{itemize}

\section{Access Requirements}

\begin{itemize}
    \item SSH access to all application nodes.
    \item Kubernetes cluster admin or namespace admin role.
    \item Access to the container registry (\texttt{ghcr.io/datapipe}).
    \item Database admin credentials for PostgreSQL.
    \item Kafka admin credentials for topic management.
\end{itemize}

%=============================================================================
\chapter{Infrastructure Provisioning}
\label{sec:provisioning}
%=============================================================================

\section{Terragrunt Configuration}

The infrastructure is managed via Terraform modules:

\begin{lstlisting}[language=bash, caption={Terraform init and apply.}]
cd deploy/terraform/production
terraform init
terraform plan -out=tfplan
terraform apply tfplan
\end{lstlisting}

\section{Kubernetes Cluster Setup}

\begin{lstlisting}[language=bash, caption={Kubernetes namespace and secrets setup.}]
# Create namespace
kubectl create namespace datapipe

# Create secrets for database credentials
kubectl create secret generic datapipe-db-credentials \
    --from-literal=username=datapipe \\
    --from-literal=password=$(openssl rand -base64 32) \\
    -n datapipe

# Create secrets for Kafka credentials
kubectl create secret generic datapipe-kafka-credentials \\
    --from-literal=sasl-username=datapipe \\
    --from-literal=sasl-password=$(openssl rand -base64 32) \\
    -n datapipe
\end{lstlisting}

%=============================================================================
\chapter{Deployment Steps}
\label{sec:deployment}
%=============================================================================

\section{Step 1: Deploy Kafka}

\begin{lstlisting}[language=bash, caption={Deploy Kafka using Strimzi operator.}]
helm repo add strimzi https://strimzi.io/charts/
helm repo update

helm install kafka strimzi/strimzi-kafka \\
    --namespace datapipe \\
    --values deploy/kafka/values-production.yaml \\
    --wait
\end{lstlisting}

\section{Step 2: Deploy PostgreSQL}

\begin{lstlisting}[language=bash, caption={Deploy PostgreSQL with Helm.}]
helm repo add bitnami https://charts.bitnami.com/bitnami

helm install postgresql bitnami/postgresql \\
    --namespace datapipe \\
    --values deploy/postgresql/values-production.yaml \\
    --set auth.postgresPassword=$DB_PASSWORD \\
    --set primary.persistence.size=100Gi \\
    --wait
\end{lstlisting}

\section{Step 3: Deploy DataPipe}

\begin{lstlisting}[language=bash, caption={Deploy DataPipe application.}]
helm install datapipe deploy/datapipe \\
    --namespace datapipe \\
    --values deploy/datapipe/values-production.yaml \\
    --set image.tag=3.2.1 \\
    --set replicas=3 \\
    --set resources.requests.cpu=8 \\
    --set resources.requests.memory=32Gi \\
    --set resources.limits.cpu=16 \\
    --set resources.limits.memory=64Gi \\
    --wait
\end{lstlisting}

\section{Step 4: Verify Deployment}

\begin{lstlisting}[language=bash, caption={Verify all pods are running.}]
# Check pod status
kubectl get pods -n datapipe

# Check logs for errors
kubectl logs -l app=datapipe -n datapipe --tail=100

# Run health check
curl -s https://datapipe.internal/api/v1/cluster/health | jq .
\end{lstlisting}

\section{Step 5: Configure Pipelines}

\begin{lstlisting}[language=bash, caption={Register production pipelines.}]
# Apply pipeline configurations
kubectl apply -f deploy/pipelines/ -n datapipe

# Verify pipeline registration
kubectl exec -it datapipe-0 -n datapipe -- \\
    dp pipeline list
\end{lstlisting}

%=============================================================================
\chapter{Health Checks}
\label{sec:health-checks}
%=============================================================================

\section{Application Health}

\begin{lstlisting}[language=bash, caption={Application health checks.}]
# Basic health check
curl -s https://datapipe.internal/api/v1/health

# Detailed cluster status
curl -s https://datapipe.internal/api/v1/cluster/status | jq .

# Pipeline status
curl -s https://datapipe.internal/api/v1/pipelines | jq .
\end{lstlisting}

\section{Infrastructure Health}

\begin{lstlisting}[language=bash, caption={Infrastructure health verification.}]
# Kubernetes node status
kubectl get nodes -o wide

# Pod resource usage
kubectl top pods -n datapipe

# Kafka cluster status
kubectl exec kafka-0 -n datapipe -- \\
    kafka-broker-api-versions.sh \\
    --bootstrap-server localhost:9092

# PostgreSQL replication status
kubectl exec postgresql-primary-0 -n datapipe -- \\
    psql -U postgres -c "SELECT * FROM pg_stat_replication;"
\end{lstlisting}

%=============================================================================
\chapter{Rollback Procedure}
\label{sec:rollback}
%=============================================================================

\section{Rollback Triggers}

Initiate a rollback if any of the following occur within 30 minutes of
deployment:

\begin{itemize}
    \item Error rate exceeds 5\% of total requests.
    \item Pipeline throughput drops below 50\% of baseline.
    \item P99 latency exceeds 5 seconds.
    \item Data loss or corruption is detected.
\end{itemize}

\section{Rollback Steps}

\begin{lstlisting}[language=bash, caption={Helm rollback procedure.}]
# List previous releases
helm history datapipe -n datapipe

# Rollback to previous version
helm rollback datapipe <REVISION> -n datapipe

# Verify rollback
kubectl get pods -n datapipe
curl -s https://datapipe.internal/api/v1/cluster/health | jq .
\end{lstlisting}

\section{Database Rollback}

If the deployment includes database migrations:

\begin{lstlisting}[language=bash, caption={Database migration rollback.}]
# List pending migrations
kubectl exec datapipe-0 -n datapipe -- dp db migrate status

# Rollback last migration
kubectl exec datapipe-0 -n datapipe -- dp db migrate rollback

# Verify database state
kubectl exec postgresql-primary-0 -n datapipe -- \\
    psql -U datapipe -c "SELECT version FROM schema_migrations;"
\end{lstlisting}

%=============================================================================
\chapter{Post-Deployment Verification}
\label{sec:post-deployment}
%=============================================================================

After deployment or rollback, perform the following checks:

\begin{enumerate}
    \item Verify all pods are in \texttt{Running} state.
    \item Confirm health check endpoint returns \texttt{200 OK}.
    \item Check pipeline status for all production pipelines.
    \item Monitor error rates for 30 minutes.
    \item Verify consumer lag is within acceptable bounds.
    \item Run a sample pipeline execution end-to-end.
    \item Update the deployment log with timestamp and version.
\end{enumerate}

\section{Deployment Log}

\begin{table}[htbp]
    \centering
    \caption{Deployment log template.}
    \label{tab:deploy-log}
    \begin{tabular}{@{}llll@{}}
        \toprule
        Date & Version & Deployer & Notes \\
        \midrule
        2025-06-15 & 3.2.0 & @alex-feng & Initial deployment \\
        2025-06-16 & 3.2.1 & @jordan-lee & Hotfix for memory leak \\
        \bottomrule
    \end{tabular}
\end{table}

\printbibliography[heading=bibintoc, title={References}]

\end{document}
